
\documentclass[unicode,10pt,a4paper,oneside,numbers=endperiod,openany]{scrartcl}



\usepackage{ifthen}
\usepackage[utf8]{inputenc}
\usepackage{graphics}
\usepackage{graphicx}
\usepackage{textcomp}
\usepackage{tcolorbox}
\usepackage{hyperref}
\usepackage{svg}
\usepackage[section]{placeins}
\usepackage[backend=biber,style=numeric]{biblatex}
\addbibresource{sources.bib}

\pagestyle{plain}
\voffset -5mm
\oddsidemargin  0mm
\evensidemargin -11mm
\marginparwidth 2cm
\marginparsep 0pt
\topmargin 0mm
\headheight 0pt
\headsep 0pt
\topskip 0pt
\textheight 255mm
\textwidth 165mm

\newcommand{\duedate} {}
\newcommand{\setduedate}[1]{%
\renewcommand\duedate {Due date:~ #1}}
\newcommand\isassignment {false}
\newcommand{\setassignment}{\renewcommand\isassignment {true}}
\newcommand{\ifassignment}[1]{\ifthenelse{\boolean{\isassignment}}{#1}{}}
\newcommand{\ifnotassignment}[1]{\ifthenelse{\boolean{\isassignment}}{}{#1}}

\newcommand{\assignmentpolicy}{
\begin{tcolorbox}[colback=gray!10, colframe=black,
                  title={HPC Lab for CSE 2025 --- Submission Instructions},
                  fonttitle=\bfseries,
                  sharp corners,
                  boxrule=0.8pt,
                  left=5pt, right=5pt, top=5pt, bottom=5pt]
% \small
\textbf{Following these instructions is mandatory}: Submissions that do not
comply \textbf{will not be considered}.
\begin{itemize}
\item \textbf{Submission Platform}:
      Assignments must be submitted via
      \href{https://sam-up.math.ethz.ch/?lecture=401-3670-00}{SAM-UP}.
\item \textbf{Required Files}: Your submission must include:
      \begin{itemize}
      \item \textbf{Report}: A PDF summarizing your results and observations,
            based on this \LaTeX\ template.
            \begin{itemize}
            \item File name:
                  \texttt{project\_number\_lastname\_firstname.pdf}
            \item You are allowed to discuss all questions with anyone you like;
                  however: (i) your submission must list anyone you discussed
                  problems with and (ii) you must write up your submission
                  independently.
            \end{itemize}
      \item \textbf{Source Code Archive}: A \textbf{compressed tar archive},
            named \texttt{project\_number\_lastname\_firstname.tar.gz},
            containing:
            \begin{itemize}
            \item All source code files.
            \item Build files and scripts (e.g., \texttt{Makefile}, batch job
                  scripts, plot scripts, etc).
                  If you modified any provided files, ensure they still work as
                  expected.
            \item The full \LaTeX\ source of your report including all needed
                  files to build the report (figures, listings, etc).
            \end{itemize}
      \item \textbf{Size Limit}: The total submission (PDF + archive) must not
            exceed \textbf{25 MB}.
      \item \textbf{File Naming Rules}: Use only \texttt{ASCII} characters
            (avoid spaces, special symbols, and non-English letters).
      \end{itemize}
\item \textbf{Grading Process}: The TAs will evaluate your project based on:
      \begin{itemize}
      \item Your report write-up.
      \item Your code implementation.
      \item Benchmarking your code's performance.
      \end{itemize}
\end{itemize}
\end{tcolorbox}
}
\newcommand{\punkte}[1]{\hspace{1ex}\emph{\mdseries\hfill(#1~\ifcase#1{Points}\or{Points}\else{Points}\fi)}}


\newcommand\serieheader[6]{
\thispagestyle{empty}%
\begin{flushleft}
\includegraphics[width=0.4\textwidth]{ETHlogo_13}
\end{flushleft}
  \noindent%
  {\large\ignorespaces{\textbf{#1}}\hspace{\fill}\ignorespaces{ \textbf{#2}}}\\ \\%
  {\large\ignorespaces #3 \hspace{\fill}\ignorespaces #4}\\
  \noindent%
  \bigskip
  \hrule\par\bigskip\noindent%
  \bigskip {\ignorespaces {\Large{\textbf{#5}}}
  \hspace{\fill}\ignorespaces \large \ifthenelse{\boolean{\isassignment}}{\duedate}{#6}}
  \hrule\par\bigskip\noindent%  \linebreak
 }

\makeatletter
\def\enumerateMod{\ifnum \@enumdepth >3 \@toodeep\else
      \advance\@enumdepth \@ne
      \edef\@enumctr{enum\romannumeral\the\@enumdepth}\list
      {\csname label\@enumctr\endcsname}{\usecounter
        {\@enumctr}%%%? the following differs from "enumerate"
        \topsep0pt%
        \partopsep0pt%
        \itemsep0pt%
        \def\makelabel##1{\hss\llap{##1}}}\fi}
\let\endenumerateMod =\endlist
\makeatother


\usepackage{listings}
\usepackage{xcolor}

\definecolor{codegreen}{rgb}{0,0.6,0}
\definecolor{codegray}{rgb}{0.5,0.5,0.5}
\definecolor{codepurple}{rgb}{0.58,0,0.82}
\definecolor{backcolour}{rgb}{0.95,0.95,0.95}

\lstdefinestyle{mystyle}{
    backgroundcolor=\color{backcolour},
    commentstyle=\color{codegreen},
    keywordstyle=\color{magenta},
    numberstyle=\tiny\color{codegray},
    stringstyle=\color{codepurple},
    basicstyle=\ttfamily\footnotesize,
    breakatwhitespace=false,
    breaklines=true,
    captionpos=b,
    keepspaces=true,
    numbers=left,
    numbersep=5pt,
    showspaces=false,
    showstringspaces=false,
    showtabs=false,
    tabsize=2
}

\lstset{style=mystyle}


\begin{document}


\setassignment
\setduedate{09 March 2026, 18:00~(Europe/Zurich)}

\serieheader{High-Performance Computing Lab for CSE}{2026}
            {Student: Benjamin Diethelm}
            {Discussed with: Lukas Ammann \linebreak \null\hfill Raphael Caixeta \linebreak \null\hfill Gian Laager}
            {Report for Project 1}{}
\newline

% you may comment this for the final submission, but make sure to comply with
% the policy
\assignmentpolicy

\section{Euler warm-up [10 points]}

\begin{enumerate}
    \item This tool is used to dynamically load packages as required, minimizing version conflicts and ensuring only the necessary components are loaded. It also sets environment variables as needed. Modules can be loaded or unloaded, for example:
    \begin{verbatim}
    module load stack/2025-06
    module unload stack/2025-06
    \end{verbatim}
    or any other module available on the system. Additionally, the commands
    \begin{verbatim}
    module avail
    module spider <module name>
    \end{verbatim}
    can be used to display a list of available modules or to get more information on a specific module respectively.

    \item This tool is employed to manage jobs on Euler or other supercomputers. Slurm is accessed via the login nodes and facilitates access to the compute nodes. Users can request specific hardware resources and submit code for execution. Slurm verifies user permissions for the requested hardware and queues the job. Once the hardware becomes available, the code is executed on the compute node.

    \item The C++ code in \texttt{01/hello\_world.cpp} outputs \texttt{Hello World from \{\$HOSTNAME\}}, where \texttt{\{\$HOSTNAME\}} represents the system's hostname.

    \item The script is located in \texttt{01/job\_1\_4.sh}, with its output and error logs in \texttt{01/1\_4.out} and \texttt{01/1\_4.err}.
    \item The script is located in \texttt{01/job\_1\_5.sh}, with its output and error logs in \texttt{01/1\_5.out} and \texttt{01/1\_5.err}.
\end{enumerate}


\section{Performance characteristics [50 points]}

\subsection{Peak performance [16 Points]}

The AMD website (\cite{amd_epyc_7763_spec}, \cite{amd_epyc_7H12_spec}) provides specifications for this model, while the Euler CPU Node Documentation \cite{euler_docs} offers details specific to Euler VII. Reference \cite{uops} is used for FMA and superscalability factors. The base clock frequency is taken from \cite{amd_epyc_7H12_spec} and \cite{amd_epyc_7763_spec}.
\begin{itemize}
\item Epyc 7H12 / Phase 1:
$$P_\textup{Core} = 2 \times 4 \times 2 \times 2.6 \textup{GHz} = 41.6\,\mathrm{GFlop/s}$$
$$P_\textup{CPU} = 64 \times P_\textup{core} = 2662.4\,\mathrm{GFlop/s}$$
$$P_\textup{node} = 2 \times P_\textup{CPU} = 5324.8\,\mathrm{GFlop/s}$$
$$P_\textup{cluster} = 292 \times P_\textup{node} = 1554.8\,\mathrm{TFlop/s}$$

\item Epyc 7763 / Phase 2:
  $$P_\textup{Core} = 2 \times 4 \times 2 \times 2.45 \textup{GHz} = 39.2\,\mathrm{GFlop/s}$$
  $$P_\textup{CPU} = 64 \times P_\textup{core} = 2508.8\,\mathrm{GFlop/s}$$
  $$P_\textup{node} = 2 \times P_\textup{CPU} = 5017.6\,\mathrm{GFlop/s}$$
  $$P_\textup{cluster} = 248 \times P_\textup{node}  = 1244.3\,\mathrm{TFlop/s}$$

\item Total:
  $$1554.8\,\mathrm{TFlop/s} + 1244.3\,\mathrm{TFlop/s} = 2799.1\,\mathrm{TFlop/s}$$

\end{itemize}

\subsection{Memory Hierarchies [8 Points]}
\subsubsection{Cache and main memory size}
Using the commands provided in the task description, we obtain the data in \autoref{tab:perf_mem_7763}. The output files are available as \texttt{02/Memory\_7H12.txt} and \texttt{02/Memory\_7763.txt}, while the corresponding PDF outputs are included as \ref{fig:epyc_7h12} and \ref{fig:epyc_7763}. The job files used to produce this output are \texttt{02/job\_2\_1.sh} and \texttt{02/job\_2\_2.sh}.

The key difference is that the EPYC 7H12 shares L3 cache across 4 cores while the EPYC 7763 shares it across 8 cores. Main memory is partitioned into 4 NUMA nodes per CPU, yielding 8 NUMA nodes per node, for both phase 1 and 2.

\begin{table}[hbt]
     \centering
     \begin{tabular}{| l | c |}
       \hline
       Main memory \hspace{1.5cm} &  256 GB \\
       \hline
       L3 cache & 512 MiB \\
       \hline
       L2 cache & 64 MiB \\
       \hline
       L1i cache & 4 MiB \\
       \hline
       L1d cache & 4 MiB \\
       \hline
     \end{tabular}
     \caption{Memory hierarchy of an Euler VII (phase 1/2) node with an AMD EPYC 7H12/7763 CPU.
              Each core has its own L1 and L2 cache. L3 cache is shared among groups of 4 cores in phase 1 and in groups of 8 cores in phase 2.}
     \label{tab:perf_mem_7763}
 \end{table}

 \begin{figure}[h]
   \centering
   \includegraphics[width=\textwidth]{images/EPYC_7H12.pdf}
   \caption{Topology of an Euler VII (phase 1) node with an AMD EPYC 7H12 CPU.}
   \label{fig:epyc_7h12}
 \end{figure}

 \begin{figure}[h]
   \centering
   \includegraphics[width=\textwidth]{images/EPYC_7763.pdf}
   \caption{Topology of an Euler VII (phase 2) node with an AMD EPYC 7763 CPU.}
   \label{fig:epyc_7763}
 \end{figure}


\subsection{Bandwidth: STREAM benchmark [10 Points]}

The following resources are available: source code (\texttt{02/stream.c.c}), batch scripts (\texttt{02/job\_2\_3.sh},  \texttt{02/job\_2\_4.sh}), output files (\texttt{02/2\_3\_7H12.out}, \texttt{02/2\_3\_7763.out})

\autoref{tab:stream_7h12} and \autoref{tab:stream_7763} show the output from STREAM. Across all measurements, phase 2 shows better results. For both phases, Add and Triad are roughly consistent, while Scale is slightly lower. Copy throughput is notably higher for both, at approximately $31.0\ \mathrm{GB/s}$ for phase 1 and $36.6\ \mathrm{GB/s}$ for phase 2


\begin{table}[hbt]
   \centering
   \begin{tabular}{| l | c | c | c | c |}
     \hline
     Function & Best Rate MB/s & Avg time & Min time & Max time \\
     \hline
     Copy & 31020.6 & 0.150761 & 0.140294 & 0.163701 \\
     \hline
     Scale & 18001.0 & 0.260390 & 0.241764 & 0.284396 \\
     \hline
     Add & 20236.0 & 0.352662 & 0.322594 & 0.428797 \\
     \hline
     Triad & 20139.2 & 0.356431 & 0.324144 & 0.466777 \\
     \hline
   \end{tabular}
   \caption{STREAM benchmark results for AMD EPYC 7H12 (Euler VII Phase 1). AI has been used to create and format this table from the slurm output files.}
   \label{tab:stream_7h12}
\end{table}
\begin{table}[hbt]
   \centering
   \begin{tabular}{| l | c | c | c | c |}
     \hline
     Function & Best Rate MB/s & Avg time & Min time & Max time \\
     \hline
     Copy & 36563.5 & 0.131608 & 0.119026 & 0.161190 \\
     \hline
     Scale & 22009.5 & 0.218883 & 0.197733 & 0.243171 \\
     \hline
     Add & 25126.1 & 0.281478 & 0.259810 & 0.359157 \\
     \hline
     Triad & 25357.8 & 0.279042 & 0.257436 & 0.319567 \\
     \hline
   \end{tabular}
   \caption{STREAM benchmark results for AMD EPYC 7763 (Euler VII Phase 2). AI has been used to create and format this table from the slurm output files.}
   \label{tab:stream_7763}
\end{table}


\subsection{Performance model: A simple roofline model [16 Points]}

The plot was generated using \texttt{02/roofline\_plot.py}, refer to \autoref{fig:roofline_plot}. The script also calculates the minimum Operational Intensity required to achieve compute-bound performance. For Euler VII phase 1, the value is $1.34\ \mathrm{Flops}/\mathrm{Byte}$, and for phase 2, it is $1.07\ \mathrm{Flops}/\mathrm{Byte}$.

\begin{figure}[h]
   \centering
   \includegraphics[width=0.8\textwidth]{images/roofline_plot.pdf}
   \caption{Roofline model plot for Euler VII nodes. AI was used in creating this plot.}
   \label{fig:roofline_plot}
\end{figure}

\section{Auto-vectorization [10 points]}

\begin{enumerate}
    \item    Data alignment refers to storing objects at memory addresses that are multiples of their own size. Misaligned memory access can result in significant performance penalties. SIMD/SSE instructions prefer 16-byte alignment for optimal efficiency. Data alignment can be enforced using \texttt{\_\_declspec(align(n))} or \texttt{\_\_builtin\_assume\_aligned}, or it may be handled automatically by the compiler.
    \item Possible reasons include:
        \begin{itemize}
            \item Non-contiguous memory access
            \item Data-dependencies between loop iterations
            \item Pointer aliasing
            \item Non-countable loops
            \item Function calls within loops
            \item Mixing data types
        \end{itemize}
    \item Possible ways include:
        \begin{itemize}
            \item \verb|#pragma ivdep| (ignore assumed dependencies)
            \item \verb|#pragma vector always| (force vectorization)
            \item \verb|#pragma vector align| (assert the data is aligned)
            \item \verb|#pragma loop count(n)| (provide typical trip count for loop)
            \item \verb|#pragma omp simd| (force SIMD)
        \end{itemize}
    \item Strip-mining splits a loop into a vectorized part and a scalar cleanup loop. Loop blocking partitions the iteration space into smaller blocks to improve cache reuse. Loop interchange reorders nested loops to obtain stride-1 memory access patterns. Loop peeling performs a few scalar iterations until an aligned boundary is reached. Loops are unrolled automatically if possible.
    \item We can use user-mandated vectorization with \verb|#pragma omp simd|, which forces vectorization and issues a warning if it fails. For more control, SIMD-enabled functions like \verb|__attribute__((vector))| or \verb|#pragma omp declare simd| allow writing scalar-style functions that the compiler can transform into short vector variants. Vector intrinsics, e.g. \verb|_mm_add_ps| or inline assembly give full manual control.
\end{enumerate}

\pagebreak

\section{Matrix multiplication optimization [30 points]}
\subsection{dgemm\_jki.c}
\label{subsec:jki}
The loop order was modified to \verb|jki| to enhance cache locality, following standard GEMM optimization principles \cite{goto_gemm,blis_toms}. \autoref{fig:variants} compares this approach to the naive implementation, OpenBLAS, and other optimized versions.
The performance of \verb|dgemm_jki.c| demonstrates significant improvement over the naive implementation, achieving $13.21\,\mathrm{GFlop/s}$ compared to $3.41\,\mathrm{GFlop/s}$ for the naive version. This 3.9 speedup results from the JKI loop order's stride-one access pattern for matrix \verb|C| in the innermost loop, which dramatically reduces cache misses compared to the naive IJK ordering. Despite non-contiguous access for matrix \verb|A|, the JKI ordering achieves $26.3\%$ of peak performance on average versus only $4.6\%$ for the naive implementation, confirming that optimizing for cache-friendly access patterns in the innermost loop provides substantial performance benefits. The other results from \autoref{fig:variants} will be discussed in \ref{subsec:analysis}.





\begin{figure}[h]
  \centering
  \includegraphics[width=\textwidth]{code/04/plots/variants/optimization_levels_performance.pdf}
  \caption{Performance comparison of the naive, jki, OpenBLAS and 1/2/3-level blocking versions. AI was used in creating this plot.}
  \label{fig:variants}
\end{figure}

\subsection{dgemm\_blocked.c}
The final optimized file is included as \verb|dgemm_blocked.c|.

\subsubsection{System}
Unless otherwise specified, all tests were conducted on a single core of an AMD EPYC 7763. As detailed in \ref{subsec:analysis}, the GCC compiler (version 13.2.0) was used, with flags \verb|-O3 -march=znver3 -ffast-math -funroll-loops| \cite{gcc_manual}. The default optimization strategy employed was 2-level blocking, as described in \ref{subsubsec:overview}.

\subsubsection{Overview}
\label{subsubsec:overview}
The optimizations implemented are:

\begin{itemize}
  \item \textbf{Multi-level blocking for cache hierarchy}: The code implements three configurable blocking strategies (1-level, 2-level, and 3-level blocking) controlled by preprocessor directives (\verb|L3_ONLY|, \verb|L2_ONLY|, \verb|L1_ONLY|). The program defaults to 2-level blocking (\verb|L2_ONLY|). The 1-level blocking targets the L3 cache with block size \verb|S3|, 2-level blocking adds L2 cache optimization with block size \verb|S2|, and 3-level blocking further incorporates L1 cache optimization with block size \verb|S1|. This multi-level cache blocking approach is a standard technique in high-performance GEMM implementations \cite{goto_gemm,blis_toms}. All three versions utilize the same micro-kernel implementation to minimize code duplication and maintain consistency. Unless stated, we use 2-level blocking everywhere.

    \item \textbf{Hierarchical block size optimization}: A systematic tuning process was employed to find optimal block sizes for each cache level. The script \verb|04/L3_opt.sh| tests all feasible block sizes for the L3 cache (values that fit in cache and are multiples of 8 to align with cache lines). Using the optimal L3 block size, \verb|L2_opt.sh| then searches for the optimal L2 block size. For the 3-level blocking variant, \verb|04/L1L2_opt.sh| performs a search over all combinations of L2 and L3 block sizes to identify the optimal configuration. This kind of empirical parameter search is a common strategy in high-performance BLAS implementations \cite{atlas}. This is further discussed in \ref{subsec:analysis}.

    \item \textbf{Data packing for spatial locality}: The micro-kernel \verb|do_block_micro| packs both matrices A (M×K) and B (K×N) into contiguous local buffers (\verb|pA| and \verb|pB|) with optimal memory layouts. Matrix A is packed with leading dimension M, and matrix B with leading dimension K. This eliminates strided memory access patterns and ensures sequential access during the computation phase, maximizing cache line utilization \cite{goto_gemm,blis_toms}.

    \item \textbf{Register-level accumulation}: The micro-kernel loads the corresponding C tile into a local accumulator array (\verb|cij|) before computation. All updates to C are performed on this local buffer throughout the computation, significantly reducing memory writes. The accumulated results are written back to the main C matrix only once per micro-kernel invocation. Register blocking/accumulation is a standard component of high-performance GEMM micro-kernels \cite{goto_gemm,blis_toms}.

    \item \textbf{Loop dependency annotation}: The code uses \verb|#pragma GCC ivdep| (ignore vector dependencies) in all innermost loops to inform the compiler that loop iterations are independent and can be safely vectorized \cite{gcc_manual}. This pragma was chosen over \verb|#pragma omp simd| after testing with \verb|04/flag_compare.sh| demonstrated that the compiler could auto-vectorize effectively without explicit OpenMP SIMD directives.

    \item \textbf{Pointer aliasing control}: All function parameters use the \verb|restrict| keyword to guarantee to the compiler that pointers do not alias, enabling more aggressive optimizations including better vectorization and instruction reordering \cite{gcc_manual}.

    \item \textbf{Compiler flag optimization}: The script \verb|04/flag_compare.sh| compares different compiler optimization flags on the 2-level blocking implementation to quantify the individual impact of each flag and identify the optimal compilation configuration \cite{gcc_manual,atlas}. This is further discussed in \ref{subsec:analysis}.

    \item \textbf{Loop ordering optimization}: The computation loop in the micro-kernel follows a j-k-i ordering, which allows the innermost loop over i to access packed matrix A sequentially while broadcasting a single element from packed matrix B, maximizing vectorization efficiency, see \ref{subsec:jki} \cite{goto_gemm,blis_toms}.

    \item \textbf{Intel Compiler}: The Intel compiler was also evaluated. However, as its performance was significantly worse than that of GCC, it was quickly abandoned.
\end{itemize}

\subsubsection{Analysis}
\label{subsec:analysis}
This section analyzes the impact of the optimizations introduced in \verb|dgemm_blocked.c| and relates them to the overall performance comparison shown in \autoref{fig:variants}. In all plots, two complementary metrics are reported: the average percentage of peak (a good indicator of sustained performance across problem sizes) and the peak throughput in $\mathrm{GFlop/s}$ (which highlights best-case behavior under favorable cache and alignment conditions). The goal of the optimizations is not only to increase compute throughput, but also to make data movement predictable so that the implementation behaves like a compute-bound kernel over a wide range of matrix sizes. Overall, the structure of the optimized implementation (cache blocking + packing + micro-kernel accumulation) is aligned with widely used high-performance GEMM designs \cite{goto_gemm,blis_toms}.

The incremental results in \autoref{fig:steps} show that the largest single improvement comes from switching the loop order to \verb|jki|, consistent with \ref{subsec:jki}. Updating contiguous elements of C in the innermost loop improves spatial locality and makes it easier for the compiler to vectorize the reduction. Subsequent steps such as packing and register accumulation improve the sustained average further, but their effect is smaller and less uniform across sizes. This is expected: packing reduces address-generation overhead and eliminates strided accesses, but introduces extra copy work that only pays off when the packed panels are reused enough times. Register-level accumulation reduces traffic to C, but increases register pressure and can constrain instruction scheduling, which can reduce the observed peak throughput even if the average improves.

The combined comparison in \autoref{fig:variants} shows how these kernel-level improvements interact with blocking and with a tuned library implementation (OpenBLAS). The naive implementation is dominated by poor locality and repeatedly streams A and B with little reuse. The \verb|jki| variant already reaches a much higher fraction of single-core performance by fixing the most critical access pattern, but it still lacks a mechanism to keep the working set in cache as sizes grow. Multi-level blocking addresses this by structuring computation so that panels of A and B are reused while resident in cache. In practice, 2-level blocking provides the best balance on this CPU: it reduces capacity misses substantially without introducing excessive packing overhead or overly small tiles. The 3-level variant adds another tiling layer, but the extra loop overhead and more frequent packing can outweigh the benefit of tighter L1 targeting, which is consistent with its lower performance in \autoref{fig:variants}.


\begin{figure}[h]
  \centering
  \includegraphics[width=\textwidth]{code/04/plots/steps/optimization_steps_performance.pdf}
  \caption{Performance comparison of the optimizations implemented in dgemm\_blocked.c. AI was used in creating this plot.}
  \label{fig:steps}
\end{figure}

As outlined in \ref{subsubsec:overview}, the scripts \verb|04/L3_opt.sh|, \verb|04/L2_opt.sh| and \verb|04/L1L2_opt.sh| explore feasible block sizes systematically to find settings that keep the working set within a targeted cache level while still providing enough reuse to amortize packing and loop overhead.

For 1-level blocking, \autoref{fig:l3} shows that performance generally improves as the L3 block size increases (up to the point where the working set begins to stress associativity, prefetching effectiveness, or translation lookaside buffers). The presence of relatively good performance at very small block sizes is also intuitive: when the tiles fit into smaller caches, even a nominal “L3-only” strategy can benefit from L2/L1 residency.

For 2-level blocking, \autoref{fig:l2} exhibits a pronounced optimum at $S2 = 64$ (with $S3$ fixed at the best value from \autoref{fig:l3}). This sharp peak indicates that once the L2 block becomes too large, the packed panels and the active C tile no longer fit comfortably in L2, increasing conflict and capacity misses; when it is too small, the compute work per packed byte is reduced and overhead becomes dominant. The selected $S2 = 64$ therefore represents a sweet spot between locality and overhead on the EPYC 7763 core.

For 3-level blocking, the slice plot and heatmap (\autoref{fig:l1l2_slices} and \autoref{fig:l1l2_heatmap}) confirm that only a narrow region of $(S1,S2)$ combinations performs well, with the best result at $S1 = 32$ and $S2 = 64$. The overall lower performance compared to 2-level blocking suggests that the additional L1-level partitioning constrains the micro-kernel too tightly and increases the frequency of packing and boundary handling. In other words, the smaller tiles improve locality but can reduce instruction-level efficiency and increase overhead enough to dominate.

As mentioned in \ref{subsubsec:overview}, various compiler flags and combinations have been tested. \autoref{fig:flags} suggests that \verb|-O3 -march=znver3 -ffast-math -funroll-loops| yields the best result for this specific benchmark setup. At the same time, the plot shows that most flag combinations are clustered relatively closely, indicating that once the implementation is cache-friendly and vectorizable, the dominant factor is the algorithmic structure rather than any single “magic” compiler switch. In particular, \verb|-O3 -march=znver3| already enables the key transformations (vectorization, scheduling, and use of ISA features), while additional flags mostly trade off minor improvements against potential downsides such as increased code size (unrolling) or numerically more aggressive transformations (fast-math). The small spread in \autoref{fig:flags} also explains why rankings can change between runs: when differences are small, measurement noise (frequency variation, OS jitter, and warm cache effects) becomes visible.

Overall, the plots tell a consistent story: \autoref{fig:steps} attributes the largest single gain to choosing a cache- and SIMD-friendly loop order; \autoref{fig:l3} and \autoref{fig:l2} show that performance is highly sensitive to block sizes once blocking is introduced; \autoref{fig:l1l2_slices} and \autoref{fig:l1l2_heatmap} indicate that adding an extra blocking level narrows the region of good configurations and can introduce enough overhead to hurt; and \autoref{fig:flags} confirms that compiler tuning is secondary compared to data movement and locality. These effects culminate in \autoref{fig:variants}, where the tuned 2-level blocked implementation provides the strongest overall performance among the hand-written variants, while OpenBLAS remains a useful reference point for what is achievable with architecture-specific assembly micro-kernels and extensive autotuning.

\begin{figure}[h]
  \centering
  \includegraphics[width=\textwidth]{code/04/plots/flag/compilation_flags_performance.pdf}
  \caption{Performance comparison of various compilation flags and combinations for the optimized version dgemm\_blocked.c. AI was used in creating this plot.}
  \label{fig:flags}
\end{figure}


\begin{figure}[h]
  \centering
  \includegraphics[width=\textwidth]{code/04/plots/block/l3_block_size_performance.pdf}
  \caption{Average percentage of peak performance as a function of the L3 block size S3 for 1-level blocking. Among the 144 block sizes tested, performance varies between $21.87\%$ and $34.81\%$ of peak, with the best result at $S3 = 1152$ achieving $34.81\%$. Performance is generally higher with larger block sizes (L3 cache) or with very small block sizes that fit in smaller caches. AI was used in creating this plot.}
  \label{fig:l3}
\end{figure}
\begin{figure}[h]
  \centering
  \includegraphics[width=\textwidth]{code/04/plots/block/l2_block_size_performance.pdf}
  \caption{Average percentage of peak performance as a function of the L2 block size S2 for 2-level blocking, with the L3 block size fixed at the optimal $S3 = 1152$. Performance increases monotonically from $22.6\%$ at $S2 = 16$ to a clear peak of $51.4\%$ at $S2 = 64$, before dropping off for larger block sizes that no longer fit efficiently in the L2 cache. The optimal $S2 = 64$ is therefore selected for the final 2-level blocked implementation. AI was used in creating this plot.}
  \label{fig:l2}
\end{figure}
\begin{figure}[h]
  \centering
  \includegraphics[width=\textwidth]{code/04/plots/block/l1l2_block_slices_performance.pdf}
  \caption{Average percentage of peak performance as a function of L1 and L2 block sizes (S1 and S2) for 3-level blocking with fixed outermost block size. Among the 72 combinations tested, performance varies between $10.54\%$ and $25.26\%$ of peak, with the best result at $S1 = 32$ and $S2 = 64$ achieving $25.26\%$. AI was used in creating this plot.}
  \label{fig:l1l2_slices}
\end{figure}
\begin{figure}[h]
  \centering
  \includegraphics[width=\textwidth]{code/04/plots/block/l1l2_block_heatmap_performance.pdf}
  \caption{Heatmap showing the average percentage of peak performance as a function of L1 and L2 block sizes (S1 and S2) for 3-level blocking with fixed outermost block size. Among the 72 combinations tested, the best performance of $25.26\%$ was achieved with $S1 = 32$ and $S2 = 64$. AI was used in creating this plot.}
  \label{fig:l1l2_heatmap}
\end{figure}

\printbibliography


\end{document}
